\chapter[Conclusão]{Conclusão}

O presente trabalho propôs o desenvolvimento e a análise comparativa de duas arquiteturas de Redes Neurais Convolucionais (CNN1D e CNN2D) aplicadas à 
classificação de arritmias cardíacas, utilizando a base de dados MIT-BIH. Diferentemente de abordagens que focam exclusivamente na otimização de 
hiperparâmetros, este estudo explorou o impacto da representação dos dados na entrada da rede, contrapondo o uso de sinais brutos no domínio do 
tempo (1D) à utilização de espectrogramas no domínio da frequência (2D). Essa perspectiva permitiu avaliar não apenas a capacidade de generalização 
dos modelos, mas também o custo-benefício computacional associado ao pré-processamento.

Analisando os resultados obtidos, pode-se concluir que ambas as estratégias mostraram-se eficazes e cumpriram o objetivo de auxiliar no diagnóstico 
automatizado, atingindo acurácias superiores a 0.97. Contudo, a CNN1D obteve resultados ligeiramente melhores do que a abordagem bidimensional na média 
global. Esse desempenho superior, aliado à ausência da necessidade de conversão para imagens (STFT), sugere que a arquitetura unidimensional é mais 
viável para cenários onde a eficiência computacional é prioritária.

Apesar dos resultados satisfatórios, este trabalho apresenta algumas limitações que devem ser consideradas. A definição dos hiperparâmetros das redes 
foi realizada com base em experimentações preliminares, não sendo conduzida uma busca sistemática que pudesse garantir a obtenção da configuração ótima. 
Além disso, os resultados apresentados correspondem a uma única execução para cada modelo, não sendo realizada uma análise estatística baseada em 
múltiplas execuções independentes, o que limita a avaliação da variabilidade dos resultados. 

Outro ponto relevante refere-se ao pré-processamento dos dados, no qual foram utilizados segmentos de tamanho fixo e técnicas específicas de aumento de 
dados, cujas escolhas podem influenciar diretamente o desempenho dos modelos. No caso da abordagem 2D, a conversão dos sinais para espectrogramas envolve 
parâmetros que também impactam a qualidade da representação. Por fim, destaca-se que a avaliação foi realizada exclusivamente na base MIT-BIH, não sendo 
investigada a capacidade de generalização dos modelos em outras bases de dados.

\section{Trabalhos Futuros}

Como desdobramento deste estudo e visando o aprimoramento contínuo dos sistemas de classificação de arritmias, é sugerido as seguintes linhas de 
investigação para trabalhos futuros:

\begin{itemize}
    \item \textbf{Balanceamento de Dados:} Investigar técnicas avançadas de aumento de dados (\textit{Data Augmentation}), como o uso de Redes 
    Adversariais Generativas (GANs), para gerar amostras sintéticas das classes minoritárias (S e F), visando reduzir o viés do classificador em 
    direção à classe normal.
    
    \item \textbf{Arquiteturas Híbridas:} Explorar a combinação de camadas convolucionais (CNN) para extração espacial de características com camadas 
    recorrentes (LSTM ou GRU) para captura de dependências temporais de longo prazo, o que pode melhorar a classificação de sequências rítmicas complexas.
    
    \item \textbf{Generalização:} Avaliar o desempenho dos modelos treinados em outras bases de dados públicas (como a do desafio \textit{PhysioNet/CinC} 
    ou a base da INCART), para verificar a capacidade de generalização dos algoritmos frente a diferentes equipamentos de aquisição e perfis de pacientes.
    
    \item \textbf{Implementação em Hardware:} Portar o modelo CNN1D para sistemas embarcados (como FPGAs ou microcontroladores ARM Cortex-M), 
    avaliando o consumo energético e o tempo de inferência para aplicações em dispositivos vestíveis (\textit{wearables}).

\end{itemize}

A intenção é que esta pesquisa contribua para a evolução de diferentes metodologias capazes de classificar arritmias cardíacas com alta precisão e 
baixo custo. Dessa maneira, tais recursos computacionais poderiam ser efetivamente integrados a dispositivos médicos e vestíveis, fornecendo suporte 
robusto aos profissionais de saúde na realização de diagnósticos precoces e assertivos.